%% ============================================================
%%  Abstract & Keywords — sn-jnl / Nature Computational Science
%%  Target length: ≤150 words (NCS strict guideline)
%%  Style: "Here we show..." declarative opening
%% ============================================================

\abstract{%
As AI models exceed accelerator memory, inference latency is governed by
hierarchical data coordination rather than computation.
Here we show that this coordination is not continuously optimisable but
exhibits intrinsic regime transitions.
Characterising any system by two dimensionless ratios---the fast-memory
residency ratio ($R_C$) and the transfer-pressure ratio ($R_B$)---we
identify three operationally distinct regimes: capacity-limited,
coordination-dominated, and I/O-limited.
Beyond capacity or bandwidth thresholds, regime-agnostic strategies fail
predictably and their rankings \emph{invert}: methods reducing latency by
24\% in favourable regimes increase it by 8--12\% once bandwidth saturates.
We establish structural lower bounds proving transitions are inevitable,
empirically calibrate boundary values ($\theta_C \approx 0.50$,
$\theta_B \approx 0.40$) across five hardware platforms, and validate
regime transitions across language and vision workloads on GPU, TPU,
and Inferentia2.
These results reframe hierarchical memory optimisation as regime-aware
design, with broad implications for memory-bound computing.
}

\keywords{hierarchical memory systems, memory orchestration, AI inference,
phase transitions, regime-dependent behaviour, large language models,
computational systems}
